\section{Literature Review}
\label{sec:results-lit}
\todo[inline]{The literature-review results have not yet been updated to
reflect the revised methodology and are therefore currently inconsistent with
it.}
This section reports the search and screening of Section~\ref{sec:lit-review}
and then what the screened literature says about the review question. All
counts come from the search scripts and the screening log; the corpus is listed
in Appendix~\ref{app:lit-inventory}.
\todo{The screening decisions in 01\_Literature/\allowbreak Literature-Review/\allowbreak screening.csv
were prepared by an AI assistant. Check every decision before this section
states who screened, and obtain the 145 full texts still missing (library, ACM Digital Library).}

\subsection{Search and Screening}

\paragraph{Identification.} The searches ran on 28 October 2026; the final runs
started between 14:46 and 15:10~UTC (Appendix~\ref{app:lit-strings}).
Table~\ref{tab:lit-sources} gives the hits and records of every source. The four
databases returned 2{,}583 hits over the strands, which were 2{,}421 records
once each database's own repeats were merged; IEEE Xplore contributed most of
them (1{,}240), and its S1 query alone returned 1{,}004 hits. The eight venues
were listed completely: 25{,}196 papers, from 372 at SOUPS to 10{,}718 at CHI.
The venue rules selected 85 of them. The German PubMed strand (S5) returned no
record. Merging the database and venue records removed 111 duplicates across
sources and left 2{,}395 records from the searches. The other routes added 585
records that no search had found: 483 from citation searching alone, 64 from the
Methods Digest alone and 38 from the thesis citations alone or from more than
one of these routes.
In total, 2{,}980 records were screened. For 105 of them no abstract was
available, and 371 were screened with an abstract looked up by DOI.

\begin{table}[tbp]
\centering
\begin{singlespace}
\caption[Sources of the literature search and their yield]{Sources of the
literature search and their yield. Hits are counted per strand and include
repeats; records are counted after merging repeats within the source.
\emph{Only here}: records no other source found. \emph{Passed}: records that
passed title-and-abstract screening (Table~\ref{tab:lit-screening}).
$^a$~Strands applied locally. $^b$~Titles only.}
\label{tab:lit-sources}
\footnotesize
\setlength{\tabcolsep}{3.5pt}
\begin{tabularx}{\textwidth}{@{}>{\raggedright\arraybackslash}p{2.2cm}l>{\raggedright\arraybackslash}p{1.9cm}>{\raggedright\arraybackslash}Xrrr@{}}
\toprule
\textbf{Source} & \textbf{Kind} & \textbf{Fields} & \textbf{Hits} & \textbf{Records} & \textbf{Only here} & \textbf{Passed} \\
\midrule
\input{tables/literature-sources}% the generated rows end with \bottomrule
\end{tabularx}
\end{singlespace}
\end{table}

Figure~\ref{fig:lit-sources} shows how much each source contributed beyond the
others. Most records of the three large databases were found by that database
alone:
1{,}154 of IEEE Xplore's 1{,}240, 850 of PubMed's 911 and 173 of arXiv's 253.
The overlap is larger where it matters for screening: of the studies that
form the review evidence after full text, 140 studies, 65 were found on arXiv,
40 in PubMed, 21 by citation searching and 15 in IEEE Xplore, and 36 are in the
Methods Digest. Citation searching found the most records that passed
title-and-abstract screening (90), followed by IEEE Xplore (81), arXiv (65) and
PubMed (60). The venue lists
added few records, because the databases and citation searching had already
found most of their papers: CHI contributed 11 records that passed, PoPETs 10,
USENIX Security, SOUPS and IEEE S\&P together 7, CSCW 1, and CCS and NDSS none.

\begin{figure}[tbp]
\centering
\input{figures/literature-sources}
\caption[Records contributed by each source]{Records contributed by each
source, after merging repeats within the source (n~=~2{,}980 records; a record
found by several sources counts for each). Dark: found by this source only;
light: also found by another source. Numbers at the bars: records that passed
title-and-abstract screening (n~=~292).}
\Description{Horizontal stacked bars for 15 sources. IEEE Xplore and PubMed
have the longest bars, mostly unique records; citation searching follows.
Venue bars are short.}
\label{fig:lit-sources}
\end{figure}

\paragraph{Screening.} Table~\ref{tab:lit-screening} and
Figure~\ref{fig:lit-exclusions} give the screening by rule, and
Figure~\ref{fig:lit-flow} the whole flow. Each excluded record counts once,
under the first rule that applied, so each rule's denominator is the number of
records left after the rules before it. R3 excluded the most records (1{,}243
of 2{,}502, 49.7\,\%), for example security systems for hospitals, health
records or the Internet of Things in general. R4 removed 585 of 1{,}259 (46.5\,\%) studies of
health apps or wearables without a privacy question, R5 276 of 674 (40.9\,\%)
privacy-preserving systems for wearables and health apps, and R6 106 of 398
(26.6\,\%) studies of privacy text that neither used language models nor
analysed the wording of risk. R1 removed 423 records (14.2\,\%) that were not
research articles, among them 30 chapters and front or back matter of one
edited volume and eight retracted papers or retraction notices; R2 removed 55
works published before 2014. After R1--R6, 292 records passed to full text.

\begin{table}[tbp]
\centering
\begin{singlespace}
\caption[Screening of the literature review by rule]{Screening of the
literature review by rule, in the order applied. Each excluded record has one
primary reason. \emph{Assessed}: records that reached the rule; \emph{\%}:
excluded as a share of assessed. Records awaiting their full text leave the
count before R7 without being excluded.}
\label{tab:lit-screening}
\small
\begin{tabularx}{\textwidth}{@{}l>{\raggedright\arraybackslash}Xrrrr@{}}
\toprule
\textbf{Rule} & \textbf{Excludes} & \textbf{Assessed} & \textbf{Excluded} & \textbf{Remaining} & \textbf{\%} \\
\midrule
\input{tables/literature-screening}% the generated rows end with \bottomrule
\end{tabularx}
\end{singlespace}
\end{table}

\begin{figure}[tbp]
\centering
\input{figures/literature-exclusions}
\caption[Exclusions by rule]{Records excluded by each rule, with the number
the rule assessed and the share excluded (n~=~2{,}980 screened; 2{,}695
excluded in total). R1--R6: title and abstract; R7--R8: full text (147
assessed).}
\Description{Horizontal bars for rules R1 to R8. R3 is the longest bar
(1,243), followed by R4 (585), R1 (423), R5 (276), R6 (106), R2 (55), R8 (7)
and R7 (0).}
\label{fig:lit-exclusions}
\end{figure}

\begin{figure}[tbp]
\centering
\resizebox{\textwidth}{!}{\input{figures/literature-flow}}
\caption[Flow of records through the literature review]{Flow of records
through the literature review, after PRISMA~2020~\parencite{page2021prisma}.
$^{*}$~Records found by no search; several routes can share a record.
$^{\dagger}$~A study can answer several sub-questions. The dashed arrow marks
the 32 cited works that the rules excluded; they form the contextual
collection.}
\Description{Flow diagram from identification (database searches 2,421 and venue
searches 85 records) through the removal of 111 duplicates to 2,980 screened
records, 2,688 exclusions at title and abstract, 292 records passed to full
text, of which 145 await their full text and 147 were assessed; 140 studies are
review evidence. A side box lists the contextual collection of 32 cited
works.}
\label{fig:lit-flow}
\end{figure}

\paragraph{Full texts.} Of the 292 records, 48 were held by the project as
PDFs. For the other 244, open-access copies were sought on 28 October 2026, in
this order: arXiv, Europe PMC, usenix.org, the PoPETs site and the open-access
locations listed by OpenAlex. They yielded 100 files: 52 from arXiv, 30 article
texts from Europe PMC, 15 from other open-access locations, 2 from usenix.org
and 1 from the PoPETs site. One file was a journal issue without the article.
In all, 147 full texts were assessed. Seven failed R8: two are not research
articles in the sense of C1 (a master's thesis and an extended abstract), one
is a two-page paper, and four do not answer a sub-question on reading (a study
of location-based services, interviews on period tracking without a privacy
analysis, a classifier that is not a large language model, and a discussion of
language models for assessing policies without an
evaluation~\parencite{aydin2024assessing}). The other 140 are the review
evidence: 67 answer LQ1, 52 LQ2, 36 LQ3 and 15 LQ4. For 145 records no full text
was obtained: 110 have no open-access location at all and 34 only locations
that returned no PDF to the script, 12 of them in the ACM Digital Library; 69 of
the 145 come from citation searching, 66 from IEEE Xplore and 20 from PubMed.
R7 has not been applied to them, because the university library was not tried.
They remain in the corpus as \emph{awaiting full text}, with the sub-questions
assigned at title and abstract (102 LQ1, 38 LQ2, 13 LQ3, 12 LQ4), and no
statement below rests on them alone.

\paragraph{Existing literature.} All 43 works cited in the thesis before this
review entered the corpus. Sixteen meet the criteria and are review evidence.
The other 27 are in the contextual collection: 10 are not research articles
(R1; web pages, official statistics, a book, a book chapter and software), 5
predate 2014 (R2; among them Solove's taxonomy and Calo's boundaries of privacy
harm), 7 lie outside the domain (R3; among them Citron and Solove's typology,
the harm model of Knabenhans et al.\ and the work on annotation with language
models), and 5 analyse privacy text without language models or risk wording
(R6; OPP-115, Polisis, MAPS, the Data Safety study of Arkalakis et al.\ and the
analysis of Apple's privacy labels by Li et~al.). Five works cited since join
them: the four reporting and review guidelines of
Section~\ref{sec:lit-review} and the paper excluded at full text. Of the 116 Methods Digest
papers, 36 are review evidence, 19 are contextual (cited in the thesis), 60
were excluded and 1 awaits its full text. Citation searching from the four seed
papers yielded 529 DOIs, of which 4 had no metadata; 81 of its 525 records
passed screening, 61 of them found by no other route.

\subsection{What the Literature Shows}

\paragraph{Years and venues.} Figure~\ref{fig:lit-years} shows that the field
is young: of the 285 records that passed screening and were not excluded at full
text, 143 (50\,\%) were published from 2023 onwards, and although S3 searched
from 2022, the first LQ3 records date from 2023. The records come from a wide
range of venues
(Figure~\ref{fig:lit-venues}): 57 from health or medical journals, 49 from
security and privacy venues, 28 from HCI venues and 38 still only as
preprints. Of the 140 studies of the review evidence, 37 appeared in health or
medical journals, 24 at security and privacy venues and 10 at HCI venues; 38
are preprints, because arXiv was the main open-access source of full texts.

\begin{figure}[tbp]
\centering
\input{figures/literature-years}
\par\smallskip
{\scriptsize\setlength{\tabcolsep}{3.1pt}
\begin{tabular}{@{}l*{13}{r}@{}}
Year & 2014 & 2015 & 2016 & 2017 & 2018 & 2019 & 2020 & 2021 & 2022 & 2023 & 2024 & 2025 & 2026 \\
Review evidence & \input{tables/literature-years-review}
\end{tabular}}
\caption[Publication years of the screened literature]{Publication years of the
records that passed title-and-abstract screening and were not excluded at full
text (n~=~285: 140 review evidence and 145 awaiting full text), by first
sub-question. The row below gives the 140 studies of the review evidence. 2026
covers January to 28 October.}
\Description{Stacked bars per year from 2014 to 2026. Counts rise from 4 in 2014
to 41 in 2024 and 44 in 2025; LQ1 dominates every year, LQ3 appears from 2023.}
\label{fig:lit-years}
\end{figure}

\begin{figure}[tbp]
\centering
\input{figures/literature-venues}
\caption[Venue types of the screened literature]{Venue types of the records that
passed screening and were not excluded at full text (n~=~285), coded from the
venue name with fixed keyword rules (\texttt{extraction\_rules} in the search
configuration). Dark: review evidence (n~=~140); light: awaiting full text
(n~=~145).}
\Description{Horizontal stacked bars for seven venue types; other computing
venues, health journals and security venues are the largest groups.}
\label{fig:lit-venues}
\end{figure}

\paragraph{Evidence map.} Figure~\ref{fig:lit-map-surfaces} crosses the app
domain with the material each study analysed, coded from titles and abstracts
with fixed keyword rules for all 285 records. Wearables in general (89) and
activity and fitness (83) dominate; reproductive health follows with 42,
mental health with 23, and sleep and diet with 9 and 10. Privacy policies are
the most analysed surface (136 records), followed by traffic or code (67) and
interfaces (53). Store listings appear in two records only, one on mental
health apps and one on medical apps, and privacy labels in ten. Of the 90
records answering LQ2, 44 analyse one surface only. The communication of health
tracking apps is thus studied mostly through one document, the policy, and its
relation to app behaviour; the store listing, which every user sees before
installing, is almost unstudied.

\begin{figure}[tbp]
\centering
\input{figures/literature-map-surfaces}
\caption[Evidence map: app domain and material studied]{Evidence map of the
records that passed screening and were not excluded at full text (n~=~285):
number of records per app domain and material or surface, coded from title and
abstract with fixed keyword rules; a record can fall in several cells.
\emph{Outside health}: LQ3 and LQ4 studies of privacy text in other domains.}
\Description{Heat map with eight app domains as rows and ten surfaces as
columns. The darkest cells are privacy policies in studies outside health,
sensor data of wearables and traffic of activity apps; the store-listing column
is almost empty.}
\label{fig:lit-map-surfaces}
\end{figure}

Figure~\ref{fig:lit-map-rq} shows, for the 140 studies of the review evidence,
which part of the thesis RQ each bears on, coded by hand from the full text.
Seventy-seven name or test harms (RQ1) and 25 compare several surfaces (RQ4);
only 12 examine protection claims (RQ3) and 11 how harms are worded or
downplayed (RQ2). Only the work on privacy
washing~\parencite{brackin2026privacy} addresses both stance and protection
claims, and no study addresses all four parts. Of the eleven RQ2 studies, two
concern health: an experiment on the length and framing of privacy notices for
fitness wearables and a design study of risk communication in intimate health
technologies~\parencite{moniz2025toward}; neither analyses what providers
actually write. The closest precedent for RQ3 is a discourse analysis of
Apple's privacy text for its Health app, compared with users'
experiences~\parencite{serttas2026bounded}; it covers one app and codes the
construction of user control, not harms. Eighty studies serve as method
precedents.

\begin{figure}[tbp]
\centering
\input{figures/literature-map-rq}
\caption[Evidence map: parts of the RQ and app domain]{Evidence map of the
review evidence (n~=~140): number of studies per part of the thesis RQ and app
domain, coded by hand from the full text; a study can fall in several cells.}
\Description{Heat map with five rows (RQ1 to RQ4 and method precedent) and eight
domains. RQ1 and method precedent are filled across most domains; the RQ2 and
RQ3 rows are sparse, with most RQ2 entries outside health.}
\label{fig:lit-map-rq}
\end{figure}

\paragraph{LQ1: harms and capabilities.} The 67 LQ1 studies document
capabilities more often than harms as such: what data an adversary or a
recipient can obtain or infer, and from what. Fifty-one of them were coded with
at least one of eight capability categories, fixed before the full-text coding
(Table~\ref{tab:lit-capabilities}); the table is the input to the harm trees.
Reading stored or transmitted data through insecure Bluetooth, network or
storage practices is the most frequent (20 studies), followed by the profiling
of users by third parties that receive app data (16), often without the policy
saying
so~\parencite{huckvale2015unaddressed,huckvale2019assessment,grundy2019data,tangari2021mobile}.
Location is inferable from shared activity routes and elevation profiles, even
when users hide their homes behind the privacy zones that fitness platforms
offer~\parencite{hassan2018analysis,dhondt2022run}; personality and other
attributes are inferable from tracker data~\parencite{zufferey2023watch};
de-identified activity, gait and heart data can be re-identified; and motion
sensors of smartwatches reveal keystrokes and PINs. For reproductive health,
studies after the overturn of Roe v.~Wade add law enforcement as a recipient:
users were concerned about access by law enforcement and third
parties~\parencite{cao2024deleted}, and the policies of female health apps named
the possible disclosure of personal data to law
enforcement~\parencite{malki2024exploring}. Mental health apps showed
unnecessary permissions, insecure cryptography and leaks of personal data and
credentials~\parencite{iwaya2023privacy}. Only one study names insurers or
employers as users of the data, through participants' concerns. The 102 LQ1
records still awaiting their full text will extend the table.

\begin{table}[tbp]
\centering
\begin{singlespace}
\caption[Capabilities documented in the review evidence]{Capabilities
documented in the review evidence (LQ1, full text assessed; 51 studies coded),
with the data they need and the actor who can use them. \emph{Studies}: numbers
in Table~\ref{tab:lit-included}. Input to the reference harm trees.}
\label{tab:lit-capabilities}
\footnotesize
\begin{tabularx}{\textwidth}{@{}>{\raggedright\arraybackslash}p{3.6cm}>{\raggedright\arraybackslash}p{3.0cm}>{\raggedright\arraybackslash}p{2.4cm}r>{\raggedright\arraybackslash}X@{}}
\toprule
\textbf{Capability} & \textbf{Data needed} & \textbf{Actor} & \textbf{n} & \textbf{Studies} \\
\midrule
\input{tables/literature-capabilities}% the generated rows end with \bottomrule
\end{tabularx}
\end{singlespace}
\end{table}

\paragraph{LQ2: communication.} The 52 LQ2 studies of the review evidence
analyse policies above all: their availability, readability and
quality~\parencite{sunyaev2015availability,robillard2019availability,oloughlin2019reviewing},
their content for menstrual and female health apps~\parencite{shipp2020period,malki2024exploring}
and their match with app behaviour~\parencite{huckvale2019assessment,tangari2021mobile}.
Fewer go beyond the policy: Hutton et al.\ evaluated apps and policies against
heuristics~\parencite{hutton2018assessing}, Malki et al.\ combined policies
with an inspection of the apps~\parencite{malki2024exploring}, Hudig and
Singh studied the data flows of fertility trackers and users'
preferences~\parencite{hudig2025intimate}, and an audit of three period trackers
used them as a user would, against data-protection and intimate-privacy
criteria~\parencite{white2026experiential}. One study analysed store
descriptions together with policies, for depression
apps~\parencite{singh2023assessment}. None codes which harms the communication
names or how it qualifies them. Among the records awaiting their full text, a
discourse analysis of how period- and fertility-tracking companies describe
privacy comes closest to this thesis and must be assessed before the gap can be
stated as confirmed.

\paragraph{LQ3: language models.} The 36 LQ3 studies use large language models
to annotate, classify, extract from or assess privacy policies and terms: data
practices, contextual integrity, compliance with privacy laws or GDPR
transparency, contradictions within policies or with privacy labels, and
answers to users' questions~\parencite{rodriguez2024llm,chanenson2025automating,xie2025evaluating}.
Each reports an empirical evaluation, as LQ3 required. A few derive risk ratings
for users from policy clauses, but none codes how the provider words or
qualifies harms, which is the coding task of this thesis. Their validation
designs are the precedents for the LLM-assisted coding; the cautions of Baumann
et al.\ and Törnberg on configuration choices and validation remain in the
contextual collection~\parencite{baumann2025llm,tornberg2024best}.

\paragraph{LQ4: framing.} The 15 LQ4 studies analyse how texts present risk or
protection: vagueness and clarity of policies~\parencite{lebanoff2018automatic},
hedging and the downplaying of risk in breach
notifications~\parencite{zou2019might}, the effect of framing in privacy notices,
the communication of differential privacy, the transparency of genetic-testing
policies~\parencite{doan2024transparency}, the promises made for
privacy-enhancing techniques~\parencite{martin2025cookies}, risk communication
in intimate health technologies~\parencite{moniz2025toward}, privacy
washing~\parencite{brackin2026privacy}, disclosures hidden in regional sections
of policies, and the construction of control in Apple's Health privacy
text~\parencite{serttas2026bounded}. Only the last analyses the actual text of a
health tracking provider; the health-related framing studies are otherwise
experiments or design studies.

\paragraph{Answer to the review question and gaps.} The confirmed evidence
documents capabilities and harms for health tracking apps well, especially
insecure handling of data, third-party profiling, location and attribute
inference and re-identification, and it shows how policies of such apps have
been sampled and coded. It barely shows how providers word the harms they
acknowledge, whether they downplay them, or how they frame the protection of
PETs: one study analyses how a single provider constructs control in its privacy
text, and no study of the review evidence combines harm coverage, stance and
protection framing across several surfaces of health tracking apps. The
language-model studies validate the annotation of data practices and
compliance, not the coding of harms or stance. Three limits qualify this
answer. First, 145 records that passed title and abstract still await their
full text, most of them from IEEE Xplore and from citation searching, which
mainly found ACM papers; the review evidence therefore leans towards open-access
work (65 of the 140 studies are on arXiv). Second, the domain and surface codes
of the evidence map come from keyword rules on titles and abstracts and are
coarse. Third, the ACM Digital Library, Scopus, Web of Science and Google
Scholar were not searched, and five venues were filtered on titles only.
